% =====================================================================
%  DANH MỤC THUẬT NGỮ VÀ TỪ VIẾT TẮT
% =====================================================================
\chapter*{DANH MỤC THUẬT NGỮ VÀ TỪ VIẾT TẮT}
\addcontentsline{toc}{chapter}{DANH MỤC THUẬT NGỮ VÀ TỪ VIẾT TẮT}

{\small
\begin{longtable}{|L{3.6cm}|L{\dimexpr\textwidth-4\tabcolsep-3\arrayrulewidth-3.6cm\relax}|}
\hline
\hd{Thuật ngữ, từ viết tắt} & \hd{Giải thích}\\ \hline
\endfirsthead
\hline
\hd{Thuật ngữ, từ viết tắt} & \hd{Giải thích}\\ \hline
\endhead
API & Application Programming Interface – giao diện lập trình ứng dụng, cách ứng dụng di động trao đổi dữ liệu với máy chủ\\ \hline
APNs & Apple Push Notification service – dịch vụ thông báo đẩy của Apple cho thiết bị iOS\\ \hline
BR & Business Rule – quy tắc nghiệp vụ, được đánh mã BR01 đến BR23 (Bảng~\ref{tab:quy-tac-nghiep-vu})\\ \hline
CSDL & Cơ sở dữ liệu\\ \hline
CSV & Comma-Separated Values – định dạng tệp dữ liệu dạng bảng phân tách bằng dấu phẩy\\ \hline
FCM & Firebase Cloud Messaging – dịch vụ gửi thông báo đẩy đa nền tảng của Google\\ \hline
FEFO & First Expired, First Out – nguyên tắc “hết hạn trước, dùng trước”, áp dụng khi trừ kho thực phẩm\\ \hline
IPO & Input – Process – Output, mô hình phân tích nghiệp vụ theo đầu vào, xử lý và đầu ra\\ \hline
MSG & Mã thông điệp hệ thống hiển thị cho người dùng (Phụ lục~\ref{pl:thong-diep})\\ \hline
NFR & Non-Functional Requirement – yêu cầu phi chức năng, được đánh mã NFR01 trở đi\\ \hline
NV & Mã nhóm nghiệp vụ (NV01 đến NV08) và nghiệp vụ con (ví dụ NV02.4)\\ \hline
OTP & One-Time Password – mã xác thực dùng một lần, gửi qua thư điện tử\\ \hline
QR & Quick Response code – mã vạch hai chiều, dùng để chia sẻ mã mời vào nhóm\\ \hline
REST & Representational State Transfer – kiểu kiến trúc cho API trao đổi dữ liệu qua HTTP\\ \hline
SRS & Software Requirements Specification – đặc tả yêu cầu phần mềm\\ \hline
UC & Use case – ca sử dụng, được đánh mã UC01 đến UC53\\ \hline
UML & Unified Modeling Language – ngôn ngữ mô hình hóa thống nhất\\ \hline
Danh mục thực phẩm & Từ điển thực phẩm dùng chung do quản trị viên quản lý, gồm tên, loại, đơn vị mặc định, mã vạch và hạn dùng gợi ý\\ \hline
Hạn dùng gợi ý & Số ngày bảo quản mặc định của một thực phẩm tại một vị trí lưu trữ, dùng để đề xuất ngày hết hạn\\ \hline
Không gian dữ liệu & Phạm vi sở hữu dữ liệu; có thể là không gian cá nhân của một người dùng hoặc không gian của một nhóm gia đình\\ \hline
Lô thực phẩm & Một lần nhập một thực phẩm vào tủ lạnh, có số lượng, ngày mua, ngày hết hạn và vị trí riêng\\ \hline
Mặt hàng & Một dòng trong danh sách mua sắm, tham chiếu tới một thực phẩm trong danh mục\\ \hline
Ngưỡng nhắc N & Số ngày trước ngày hết hạn mà hệ thống bắt đầu nhắc nhở, mặc định bằng 3\\ \hline
Tỷ lệ đáp ứng R & Tỷ lệ nguyên liệu bắt buộc của một công thức đang có đủ trong tủ lạnh\\ \hline
Trưởng nhóm & Người tạo nhóm gia đình hoặc người được chuyển quyền, có quyền điều phối và phân công mua sắm\\ \hline
\end{longtable}
}
